% Part:sets-functions-relations
% Chapter: sets
% Section: schroder-bernstein

\documentclass[../../../include/open-logic-section]{subfiles}

\begin{document}

\olfileid{sfr}{siz}{sb}

\olsection{大小的概念与 Schr\"oder-Bernstein 定理}

\begin{explain}
这里有一个直观的想法：如果 $A$ 不比 $B$ 大，且 $B$ 也不比 $A$ 大，那么 $A$ 和 $B$ 就是等势的。坦率地说，如果这个想法是\emph{错的}，我们就很难再主张：我们所定义的等势概念与集合间“大小”的比较有什么关系！不过幸运的是，这个直观的想法是正确的。这一点由 Schr\"oder-Bernstein 定理所证实。
\end{explain}

\begin{thm}[Schr\"oder-Bernstein]
	\ollabel{thm:schroder-bernstein}
	如果 $\cardle{A}{B}$ 且 $\cardle{B}{A}$，
	则 $\cardeq{A}{B}$。
\end{thm}

\begin{explain}
换句话说，如果存在从 $A$ 到~$B$ 的单射，以及从 $B$ 到~$A$ 的单射，那么存在从 $A$ 到~$B$ 的双射。

然而，这个结果实际上相当\emph{难以}证明。事实上，尽管 Cantor 陈述了这一结果，但给出证明的却是他人。\footnote{关于更多历史细节，例如可参见 \citet[pp.~165--6]{Potter2004}。}
\oliflabeldef{sfr:cardinals:card-sb:sec}{我们要到 \olref[sfr][cardinals][card-sb]{sec} 节才能\emph{证明} Schr\"oder-Bernstein 定理。}{}%
目前，你只能（也必须）暂且接受它。

幸运的是，Schr\"oder-Bernstein 定理是\emph{正确的}，它证实了我们把所定义的关系——即 $\cardeq{A}{B}$ 和 $\cardle{A}{B}$——视为与“大小”有关的想法。此外，Schr\"oder-Bernstein 定理非常\emph{有用}。想出两个等势集合之间的双射往往很困难。Schr\"oder-Bernstein 定理允许我们将比较拆分开来，这样我们就只需考虑从第一个集合到第二个集合的单射，反之亦然。
\end{explain}

\end{document}